\chapter{Hereditary Angioedema}
\index{Hereditary Angioedema}
\label{sec:Hereditary Angioedema}

\section{Overview}
Hereditary angioedema (HAE) is an autosomal dominant disorder caused by deficiency (type I, 85\%) or dysfunction (type II, 15\%) of C1 esterase inhibitor (C1-INH), with a prevalence of 1 in 50,000. This genetic disorder results from specific gene mutations or chromosomal abnormalities. It may be present at birth or manifest later in life depending on the underlying mechanism. Genetic disorders result from abnormalities in an individual's DNA, ranging from single-gene mutations to chromosomal abnormalities. These conditions may be inherited from parents or arise from new mutations. Pattern of inheritance depends on whether the mutation is dominant, recessive, or X-linked. Genetic counseling helps families understand risks and make informed decisions.
\section{Causes}
This condition happens because of uncontrolled activation of the contact system, leading to excessive bradykinin production, increased vascular permeability, and recurrent episodes of subcutaneous and submucosal swelling. Genetic disorders result from mutations in specific genes that encode proteins essential for normal cellular function. The pattern of inheritance depends on whether the mutation is dominant, recessive, or X-linked.   
\section{Risk Factors}

\begin{itemize}[noitemsep]
  \item Genetic predisposition and family history
  \item Advancing age
  \item Lifestyle factors (diet, exercise, smoking, alcohol)
  \item Environmental exposures
  \item Underlying medical conditions
  \item Medication use
\end{itemize}
\section{Signs and Symptoms}

\subsection{Common symptoms}
\begin{itemize}[noitemsep]
  \item You may experience recurrent episodes of non-pitting, non-pruritic swelling affecting the extremities, face, genitals, trunk, and potentially life-threatening laryngeal swelling leading to airway obstruction. digestive tract involvement causes severe abdominal pain, nausea, vomiting, and diarrhea, often misdiagnosed as acute abdomen. Episodes last 2--5 days without treatment. Unlike allergic angioedema, there is no urticaria, itching, or response to antihistamines or epinephrine.

  \item Pain or discomfort in the affected area
  \end{itemize}

\subsection{Emergency warning signs}
\begin{itemize}[noitemsep]
  \item Severe or worsening symptoms of hereditary angioedema require immediate medical evaluation
\end{itemize}
\section{Progression and Stages}
The condition may progress through different stages of severity over time. Early stages often involve mild or subtle symptoms that may be overlooked. Moderate stages involve more noticeable symptoms affecting daily function. Advanced stages may involve significant impairment and complications.
\section{Complications}
\begin{itemize}[noitemsep]
  \item Disease progression leading to worsening symptoms
  \item Secondary infections or injuries
  \item Side effects of treatment
  \item Reduced quality of life
  \item Emotional and psychological impact
\end{itemize}
\section{Diagnosis}
Doctors diagnose this based on low C4 levels during attacks (screening test), low C1-INH antigenic levels (type I), low functional C1-INH activity (type II), and C1-INH genetic testing.

\begin{itemize}[noitemsep]
  \item Comprehensive medical history and physical examination
  \item Laboratory tests to assess organ function
  \item Imaging studies to visualize affected structures
  \item Specialized testing depending on the affected system
  \item Consultation with relevant specialists
\end{itemize}
\section{Prevention}
Treatment options include on-demand treatment for acute attacks (plasma-derived C1-INH, recombinant C1-INH, icatibant [bradykinin B2 receptor antagonist], ecallantide [kallikrein inhibitor]), short-term prophylaxis prior to procedures (plasma-derived C1-INH or attenuated androgens), and long-term prophylaxis (plasma-derived C1-INH, lanadelumab [monoclonal antibody against plasma kallikrein], berotralstat [oral plasma kallikrein inhibitor], or attenuated androgens: danazol, stanozolol).

\begin{itemize}[noitemsep]
  \item Maintain a healthy lifestyle with balanced diet and exercise
  \item Avoid known risk factors
  \item Attend regular medical check-ups for early detection
  \item Follow recommended screening guidelines
\end{itemize}
\section{Treatment Options}
Symptomatic management and supportive care Enzyme replacement therapy for certain conditions Gene therapy (emerging for select conditions) Dietary modifications (e.g., phenylketonuria) Regular monitoring for associated complications Physical and occupational therapy Genetic counseling for family planning Participation in clinical trials for new therapies

\begin{itemize}[noitemsep]
  \item Medications to manage symptoms and address underlying mechanisms
  \item Lifestyle modifications and self-care measures
  \item Physical therapy and rehabilitation
  \item Surgical intervention when indicated
  \item Regular monitoring and follow-up care
\end{itemize}
\section{Living with It}
\begin{itemize}[noitemsep]
  \item Follow your treatment plan as prescribed
  \item Maintain regular follow-up with your healthcare provider
  \item Make necessary lifestyle adjustments
  \item Seek support from family, friends, or support groups
  \item Stay informed about your condition
\end{itemize}
\section{Outlook}
\begin{itemize}[noitemsep]
  \item Most people recover well with appropriate management
  \item Before modern therapy, mortality from laryngeal swelling was 25--30\%.
\end{itemize}
